\documentclass[11pt,class=report,crop=false]{standalone}
\usepackage[screen]{../logic}


\begin{document}


%====================================================================
\chapitre{Preuves}
%====================================================================

% \insertvideo{}{partie 1.1. Titre}


\objectifs{Nous adaptons la construction des preuves aux quantificateurs de la logique du premier ordre.}


%%%%%%%%%%%%%%%%%%%%%%%%%%%%%%%%%%%%%%%%%%%%%%%%%%%%%%%%%%%%%%%%%%%%%
\section{Preuves en logique du premier ordre}

\index{preuve}

%--------------------------------------------------------------------
\subsection{Rappels}

On rappelle qu'une preuve consiste à partir d'un ensemble de formules $\mathcal{P}_1$, \ldots, $\mathcal{P}_l$ jusqu'à arriver à une formule $\mathcal{Q}$ en un nombre fini d'étapes et en suivant des règles précises. Si on a obtenu une preuve, on note :
\[
\mathcal{P}_1, \ldots, \mathcal{P}_l \quad \lturn \quad \mathcal{Q}
\]
Pour la logique du premier ordre, on reprend d'abord toutes les règles d'inférence des preuves de la logique V/F.
On renvoie au chapitre \og{}Preuves\fg{} de la première partie.
On va ajouter de nouvelles règles liées aux quantificateurs.
On rappelle qu'il n'y a encore aucun lien entre ce que l'on prouve et ce qui est vrai. 
Ce lien sera donné dans les chapitres suivants. En particulier dans le prochain chapitre on montrera que tout ce qui est prouvé est vrai.


Dans ce chapitre, on fixe un langage $\mathcal{L}$. Une preuve ne fait intervenir que ce langage. 
Il n'y a pas de structures ou d'affectations à considérer.
Ainsi les formules d'une preuve restent dans $\mathcal{L}$, ce qui est un aspect où la notion de preuve ($\lturn$)
est plus facile que la notion de satisfaction ($\ldoubleturn$).
 

%--------------------------------------------------------------------
\subsection{Règles pour les quantificateurs}

\index{regles d'inference@règles d'inférence}
\index{quantificateurs@quantificateurs $\forall$, $\exists$}

Nous présentons quatre nouvelles règles liées aux quantificateurs : il s'agit des règles d'introduction et d'élimination de $\forall$ et $\exists$.
Une règle supplémentaire concerne l'égalité.
Nous présentons chacune des règles succinctement avant d'y revenir en détails dans la section suivante.
Il y a deux règles faciles et deux règles difficiles.
La difficulté ne concerne pas la règle encadrée mais plutôt ses conditions d'application.

Pour clarifier l'écriture des règles, lorsque qu'on a une formule $\mathcal{P}$  et qu'on s'intéresse
à sa dépendance vis-à-vis de la variable $x$, alors on la note naturellement $\mathcal{P}(x)$. Attention, elle peut aussi dépendre d'autres variables.
Pour une formule $\mathcal{P}(x)$, on notera $\mathcal{P}(t)$ au lieu de $\mathcal{P}[t/x]$, la substitution de la variable $x$ par le terme $t$.



\bigskip

\textbf{Règle I. Élimination de $\forall$.} 
\mybox{$
	\forall x \ \ \mathcal{P}(x)  \quad \lturn \quad  \mathcal{P}(c)
	$}

	\begin{itemize}
	\item \emph{Conditions.} $c$ peut être une constante quelconque du langage, $c$ peut aussi être un \emph{terme clos}, c'est-à-dire un terme sans variable libre.
	
	\item \emph{Remarque.} Ainsi sous l'hypothèse $\forall x \ \ \mathcal{P}(x)$, on peut remplacer n'importe quelle occurrence libre de $x$ par $c$ et cela prouve  $\mathcal{P}(c)$ (qui est en fait $\mathcal{P}[c/x]$).
	Ici cela ne pose de problème si la constante (ou le terme) $c$ est déjà apparu auparavant dans la preuve.
	
	\item \emph{Exemple.} Pour n'importe quelle constante $c$ :
	\[
	\forall x \ \ \exists y \ \ P(x) \land Q(x,y)  \quad \lturn \quad  \exists y \ \ P(c) \land Q(c,y)
	\]
	\end{itemize}

\bigskip

\textbf{Règle II. Introduction de $\forall$.} 
\mybox{$
	\mathcal{P}(c) \quad \lturn \quad  \forall x \ \ \mathcal{P}(x)
	$}

	\begin{itemize}
	\item \emph{Conditions.} $c$ désigne une constante de $\mathcal{L}$, mais elle doit être \emph{arbitraire}, 
	c'est-à-dire qu'elle ne doit apparaître dans aucune hypothèse en cours d'utilisation. 
	
	\item \emph{Remarque.} C'est une règle piégeuse car ici $c$ ne joue pas vraiment le rôle d'une constante mais plutôt d'une \emph{variable libre}.
	Le mieux c'est d'utiliser un nom $c$ qui n'apparaît pas du tout dans les lignes précédentes de la preuve.
	
	\item \emph{Exemple.} Soit $c$ arbitraire : 
	\[
	c^2 \ge 0  \quad \lturn \quad  \forall x \ \ x^2 \ge 0
	\]
	En effet, si on a prouvé $c^2 \ge 0$ pour  un $c$ quelconque, alors pour tout $x$, $x^2 \ge 0$.
	\end{itemize}

\bigskip

\textbf{Règle III. Introduction de $\exists$.} 
\mybox{$
	\mathcal{P}(c) \quad \lturn \quad  \exists x \ \ \mathcal{P}(x)
	$}

	\begin{itemize}
	\item \emph{Conditions.} $c$ est une constante du langage ou bien un terme sans variable libre.
	
	\item \emph{Remarque.} $c$ est appelé un \emph{témoin} du quantificateur $\exists$. Ce symbole $c$ a pu apparaître auparavant dans la preuve.
	
	\item \emph{Exemple.} 
	\[
	c^2 \neq  c  \quad \lturn \quad  \exists x \ \ x^2 \neq x
	\]
	\end{itemize}
	

\bigskip
	
\textbf{Règle IV. Élimination de $\exists$.} 
\mybox{$
	\exists x \ \ \mathcal{P}(x), \ \big[\mathcal{P}(c) \lturn \mathcal{Q} \big]  \quad \lturn \quad  \mathcal{Q}
	$}
	
	La notation entre crochets correspond à une sous-preuve.
	C'est plus clair avec l'écriture sous forme d'un tableau :
	
\[
	\begin{nd}
		\hypo [~] {} {\exists x \ \ \mathcal{P}(x)}  \by{}{}
		\open				
		\hypo [~] {} {\mathcal{P}(c)} \by{\text{Sous-hypothèse}}{}
		\have [~] {} {\ \vdots} \by{}{}
		\have [~] {} {\mathcal{Q}} \by{}{}
		\close 
		\have [\lturn] {} {\mathcal{Q}} \by{}{}
	\end{nd}
\]
	
	
	\medskip
	
	\begin{itemize}
	\item \emph{Conditions.} $c$ doit être arbitraire : elle ne doit apparaître ni dans $\mathcal{Q}$, ni dans $\exists x \ \ \mathcal{P}(x)$, ni dans aucune hypothèse en cours d'utilisation.
	On parle de \emph{constante fraîche} (ou abusivement de \emph{variable fraîche}).
	
	\item \emph{Remarque.} 
	La condition n'est pas évidente au premier abord, mais le principe est naturel : si l'on sait qu'il existe un élément vérifiant $\mathcal{P}$, et que le fait de supposer $\mathcal{P}(c)$ pour un $c$ arbitraire conduit à $\mathcal{Q}$, alors la conclusion $\mathcal{Q}$ est prouvée.
	
	\item \emph{Exemple.} 
	\[
	\exists x \ \ f(x)=0, \ \forall x \ \ \big( f(x) = 0 \limplies g(x) = 1 \big) \quad \lturn \quad  \exists x \ \ g(x) = 1
	\]
	C'est une application presque directe de la règle IV, voir la section suivante pour les détails.
\end{itemize}	


\bigskip

\textbf{Règle V. Égalité.} 
\mybox{$
	t = s, \ \  \mathcal{P}(t) \quad \lturn \quad  \mathcal{P}(s)
	$}
	
	\begin{itemize}
		\item \emph{Conditions.} $t$ et $s$ sont des termes du langage. Il faut bien sûr que les substitutions par $t$ ou $s$ soient légitimes (pas de capture).
		
		\item \emph{Remarque.} Il n'y a pas de piège avec l'égalité. Deux choses égales sont considérées comme identiques et donc interchangeables à volonté.
		
		\item \emph{Exemple.}
		\[
		a=x+1, \ a > 7 \quad \lturn \quad x+1 > 7 
		\]
	\end{itemize}
	
	
	
%--------------------------------------------------------------------
\subsection{Exemples de preuves}

\index{preuve!deduction naturelle@déduction naturelle}

\begin{exemple}
	
	\[ 
	\forall x \ \ \mathcal{P}(x)  \quad \lturn \quad  \exists x \ \ \mathcal{P}(x)
	\]
	
	\medskip
	\qquad\qquad
	$
	\begin{nd}
		\hypo [~] {1} {\forall x \ \ \mathcal{P}(x)}  \by{\text{Hypothèse}}{}
		\have [~] {2} {\mathcal{P}(c)} \by{\text{Élimination du $\forall$}}{}
		\have [\lturn] {3} {\exists x \ \mathcal{P}(x)} \by{\text{Introduction du $\exists$}}{}
	\end{nd}
	$
	
	\medskip
	
	Les règles de preuve supposent implicitement que l'on travaille avec des ensembles non vides. 
	C'est cohérent avec le domaine d'une structure qui doit être un ensemble non vide.
\end{exemple}



\begin{exemple}
	Soit une formule $\mathcal{P}$ dans laquelle la variable $x$ n'apparaît pas.
	
	\[ 
	\forall x \ \ \big( \mathcal{P} \limplies \mathcal{Q}(x) \big)  
	\quad \lturn \quad  
	\mathcal{P} \limplies \big( \forall x \ \ \mathcal{Q}(x) \big)
	\]
	
	\medskip
	\qquad\qquad
	$
	\begin{nd}
		\hypo [1] {1} {\forall x \ \ \big( \mathcal{P} \limplies \mathcal{Q}(x) \big) }  \by{}{}
		\open				
			\hypo [2] {2} {\mathcal{P}} \by{\text{Sous-hypothèse}}{}
			\have [3] {3} {\mathcal{P} \limplies \mathcal{Q}(c)} \by{\text{Élimination du $\forall$ de 1}}{}
			\have [4] {4} {\mathcal{Q}(c)} \by{}{}
			\have [5] {5} {\forall x \ \ \mathcal{Q}(x)} \by{\text{Introduction du $\forall$, $c$ n'apparait dans aucune hypothèse}}{}		
		\close 
		\have [\lturn] {} {\mathcal{P} \limplies \big( \forall x \ \ \mathcal{Q}(x) \big)} \by{}{}
	\end{nd}
	$
	
\end{exemple}



\begin{exemple}

	\[ 
	\forall x \ \ \big( \mathcal{P}(x) \limplies \mathcal{Q}(x) \big), \ \   
	\exists x \ \  \mathcal{P}(x)
	\quad \lturn \quad  
	\exists x \ \  \mathcal{Q}(x)
	\]
	
	C'est assez naturel : si $\mathcal{P}(x)$ implique toujours $\mathcal{Q}(x)$ et que l'un des $\mathcal{P}(x)$ est vrai, alors l'un des $\mathcal{Q}(x)$ est vrai. C'est une variante du \emph{modus ponens} de la logique propositionnelle : $\mathcal{P}, \mathcal{P} \limplies \mathcal{Q} \ \  \lturn \ \ \mathcal{Q}$.
	
	\medskip
	\qquad\qquad
	$
	\begin{nd}
		\hypo [1] {1} {\forall x \ \ \big( \mathcal{P}(x) \limplies \mathcal{Q}(x) \big) }  \by{}{}
		\hypo [2] {2} {\exists x \ \  \mathcal{P}(x)}  \by{}{}	
		\open		
			\hypo [3] {3} {\mathcal{P}(c)} \by{\text{Sous-hypothèse}}{}
			\have [4] {4} {\mathcal{P}(c) \limplies \mathcal{Q}(c)} \by{\text{Élimination du $\forall$ de 1}}{}
			\have [5] {5} {\mathcal{Q}(c)} \by{}{}
			\have [6] {6} {\exists x \ \ \mathcal{Q}(x)} \by{\text{Introduction du $\exists$}}{}				
		\close 
		\have [\lturn] {} {\exists x \ \ \mathcal{Q}(x)} \by{\text{Élimination du $\exists$ de 2 par 3-6}}{}
	\end{nd}
	$
	
\end{exemple}



\begin{exemple}
	
	\[ 
	\exists x \ \ \forall y \ \ \mathcal{P}(x,y) 
	\quad \lturn \quad  
	\forall y \ \ \exists x \ \ \mathcal{P}(x,y) 	
	\]
	
	On peut donc intervertir un $\exists \  \forall$ vers un $\forall \ \exists$ (mais pas dans l'autre sens).
	
	\medskip
	\qquad\qquad
	$
	\begin{nd}
		\hypo [1] {1} {\exists x \ \ \forall y \ \ \mathcal{P}(x,y)}  \by{}{}
		\open		
			\hypo [2] {2} {\forall y \ \ \mathcal{P}(a,y)} \by{\text{Avec $a$ arbitraire}}{}		
			\have [3] {3} {\mathcal{P}(a,b)}{} \by{\text{Avec $b$ arbitraire}}{}
			\have [4] {4} {\exists x \ \ \mathcal{P}(x,b)} \by{\text{Introduction du $\exists$}}{}			
			\have [5] {5} {\forall y \ \ \big( \exists x \ \ \mathcal{P}(x,y) \big)} \by{\text{Introduction du $\forall$}}{}				
		\close 
		\have [\lturn] {} {\forall y \ \ \exists x \ \ \mathcal{P}(x,y)} \by{\text{Élimination du $\exists$ de 1 par 2-5}}{}
	\end{nd}
	$
	
\end{exemple}




%%%%%%%%%%%%%%%%%%%%%%%%%%%%%%%%%%%%%%%%%%%%%%%%%%%%%%%%%%%%%%%%%%%%%
\section{Négation des quantificateurs et règles dérivées}

%--------------------------------------------------------------------
\subsection{Négation des quantificateurs}

\index{quantificateurs@quantificateurs $\forall$, $\exists$}

\textbf{Règle dérivée. Négation des quantificateurs.} 
\mybox{$
	\lnot \ \big( \forall x \ \ \mathcal{P}(x) \big)  
	\quad \lturn \quad  
	\exists x \ \ \big( \lnot \mathcal{P}(x) \big) 
	$}
\mybox{$
	\lnot \ \big( \exists x \ \ \mathcal{P}(x) \big)  
	\quad \lturn \quad  
	\forall x \ \ \big( \lnot \mathcal{P}(x) \big) 
	$}

Nous ne sommes pas surpris : la négation d'un \og{}pour tout\fg{} est un \og{}il existe\fg{}, et réciproquement.

Exemples :
\begin{enumerate}
	
	\item $\lnot \ \big( \forall x \ \ f(x) = g(x) \big) \quad \lturn \quad  
	\exists x \ \ \big( f(x) \neq g(x) \big)$
	
	\item $\lnot \ \big( \exists x \ \ P(x) \land Q(x) \big) \quad \lturn \quad  
	\forall x \ \ \big( \lnot P(x) \big) \lor \big( \lnot Q(x) \big)$
	
	\item $\lnot \ \big( \forall x \ \ P(x) \limplies Q(x) \big) \quad \lturn \quad \exists x \ \ P(x) \land \big( \lnot Q(x) \big)$
\end{enumerate}


%--------------------------------------------------------------------
\subsection{Autres règles dérivées}


\textbf{Règle dérivée. Renommage.} 
\mybox{$
	\forall x \ \ \mathcal{P}(x) 
	\quad \lturn \quad  
	\forall y \ \ \mathcal{P}(y)  	
	\qquad \text{ et }\qquad
	\exists x \ \ \mathcal{P}(x) 
	\quad \lturn \quad  
	\exists y \ \ \mathcal{P}(y) 	
	$}

La variable $y$ ne doit pas apparaître comme variable libre dans la formule $\forall x \ \mathcal{P}(x)$ ni dans $\exists x \ \mathcal{P}(x)$.

\bigskip

\textbf{Règle dérivée. Interversion de deux $\forall$ ou deux $\exists$.}
\mybox{$
	\forall x \ \ \forall y \ \ \mathcal{P}(x,y)
	\quad \lturn \quad  
	\forall y \ \ \forall x \ \ \mathcal{P}(x,y)
	$}
\mybox{$
	\exists x \ \ \exists y \ \ \mathcal{P}(x,y)
	\quad \lturn \quad  
	\exists y \ \ \exists x \ \ \mathcal{P}(x,y)
	$}

\bigskip

\textbf{Règle dérivée. Interversion de $\exists \  \forall$ vers $\forall \ \exists$.}	
\mybox{
	$\exists x \ \ \forall y \ \ \mathcal{P}(x,y) 
\quad \lturn \quad  
\forall y \ \ \exists x \ \ \mathcal{P}(x,y)$	
}
On l'a prouvée dans un exemple précédent. 


\bigskip

\textbf{Règle dérivée. Commutation de $\forall$ et $\land$.}
\mybox{$
	\forall x \ \ \big( \mathcal{P}(x) \land \mathcal{Q}(x) \big)
	\quad \lturn \quad  
	\big( \forall x \ \ \mathcal{P}(x) \big) \land \big( \forall x \ \ \mathcal{Q}(x) \big)
	$}
\mybox{$
	\big( \forall x \ \ \mathcal{P}(x) \big) \land \big( \forall x \ \ \mathcal{Q}(x) \big)
	\quad \lturn \quad  
	\forall x \ \ \big( \mathcal{P}(x) \land \mathcal{Q}(x) \big)	
	$} 
Ainsi on peut passer d'un sens à l'autre. Les formules sont interdérivables.

\bigskip

\textbf{Règle dérivée. Commutation de $\exists$ et $\lor$.}
\mybox{$
	\exists x \ \ \big( \mathcal{P}(x) \lor \mathcal{Q}(x) \big)
	\quad \lturn \quad  
	\big( \exists x \ \ \mathcal{P}(x) \big) \lor \big( \exists x \ \ \mathcal{Q}(x) \big)
	$}
\mybox{$
	\big( \exists x \ \ \mathcal{P}(x) \big) \lor \big( \exists x \ \ \mathcal{Q}(x) \big)
	\quad \lturn \quad  
	\exists x \ \ \big( \mathcal{P}(x) \lor \mathcal{Q}(x) \big)	
	$} 

%--------------------------------------------------------------------
\subsection{Preuve des règles dérivées}

Voyons quelques preuves des règles dérivées.


\begin{exemple}
	Les règles de renommage sont naturelles. Les preuves sont de simples jeux d'écriture. 
	
	\[ 
	\forall x \ \ \mathcal{P}(x)  \quad \lturn \quad  \forall y \ \ \mathcal{P}(y)
	\]
	
	\medskip
	\qquad\qquad
	$
	\begin{nd}
		\hypo [~] {} {\forall x \ \ \mathcal{P}(x)}  \by{}{}
		\have [~] {} {\mathcal{P}(c)} \by{\text{Élimination du $\forall$, $c$ arbitraire}}{}
		\have [\lturn] {} {\forall y \ \mathcal{P}(y)} \by{\text{Introduction du $\forall$}}{}
	\end{nd}
	$
	
	
	\bigskip	
		
	\[ 
	\exists x \ \ \mathcal{P}(x)  \quad \lturn \quad  \exists y \ \ \mathcal{P}(y)
	\]
	\medskip
	\qquad\qquad
	$
	\begin{nd}
		\hypo [~] {} {\exists x \ \ \mathcal{P}(x)}  \by{}{}
		\open				
			\hypo [~] {} {\mathcal{P}(c)} \by{\text{$c$ constante}}{}
			\have [~] {} {\exists y \ \ \mathcal{P}(y)} \by{\text{Introduction du $\exists$}}{}
		\close 
		\have [\lturn] {} {\exists y \ \ \mathcal{P}(y)} \by{\text{Élimination du $\exists$ initial}}{}
	\end{nd}
	$	

\end{exemple}


\begin{exemple}
\[
	\forall x \ \ \forall y \ \ \mathcal{P}(x,y)
	\quad \lturn \quad  
	\forall y \ \ \forall x \ \ \mathcal{P}(x,y)	
\]

	\medskip
\qquad\qquad
$
\begin{nd}
	\hypo [~] {} {\forall x \ \ \forall y \ \ \mathcal{P}(x,y)}  \by{}{}
	\have [~] {} {\forall y \ \ \mathcal{P}(a,y)} \by{\text{Élimination du premier $\forall$, $a$ arbitraire}}{}
	\have [~] {} {\mathcal{P}(a,b)} \by{\text{Élimination du second $\forall$, $b$ arbitraire}}{}
	\have [~] {} {\forall x \ \ \mathcal{P}(x,b)} \by{\text{Introduction d'un $\forall$}}{}	
	\have [\lturn] {} {\forall y \ \ \forall x \ \ \mathcal{P}(x,y)} \by{\text{Introduction d'un autre $\forall$}}{}
\end{nd}
$
\end{exemple}


\begin{exemple}
	\[
	\big( \forall x \ \ \mathcal{P}(x) \big) \land \big( \forall x \ \ \mathcal{Q}(x) \big)
	\quad \lturn \quad  
	\forall x \ \ \big( \mathcal{P}(x) \land \mathcal{Q}(x) \big)
	\]
	
	\medskip	
	
	Voici la preuve, à vous d'expliciter les règles utilisées.
	
	\medskip
	\qquad\qquad
	$
	\begin{nd}
		\hypo [~] {} {\Big( \forall x \ \ \mathcal{P}(x) \Big) \land \Big( \forall x \ \ \mathcal{Q}(x) \Big)}  \by{}{}
		\have [~] {} {\forall x \ \ \mathcal{P}(x)} \by{}{}
		\have [~] {} {\mathcal{P}(c)} \by{}{}	
		\have [~] {} {\forall x \ \ \mathcal{Q}(x)} \by{}{}
		\have [~] {} {\mathcal{Q}(c)} \by{}{}	
		\have [~] {} {\mathcal{P}(c) \land \mathcal{Q}(c)} \by{}{}				
		\have [\lturn] {} {\forall x \ \ \big( \mathcal{P}(x) \land \mathcal{Q}(x) \big)} \by{}{}
	\end{nd}
	$
	
\end{exemple}

\begin{exemple}
	\[
	\lnot \ \Big( \forall x \ \ \mathcal{P}(x) \Big)  
	\quad \lturn \quad  
	\exists x \ \ \Big( \lnot \mathcal{P}(x) \Big) 
	\]

\medskip
\qquad\qquad
$
\begin{nd}
	\hypo [1] {} {\lnot \ \Big( \forall x \ \ \mathcal{P}(x) \Big)}  \by{}{}
	\open				
		\hypo [2] {} {\lnot \ \Big( \exists x \ \ \lnot \mathcal{P}(x) \Big)}  \by{\text{Raisonnement par l'absurde}}{}
		\open		
			\hypo [3] {} {\lnot  \mathcal{P}(c)} \by{\text{Hypothèse, avec $c$ arbitraire}}{}		
			\have [4] {} {\exists x \ \ \lnot \mathcal{P}(x)} \by{\text{Introduction de $\exists$}}{}
			\have [5] {} {\lantitauto} \by{\text{Contradiction 2 avec 4}}{}
		\close
		\have [6] {} {\mathcal{P}(c)} \by{\text{Par 2-5}}{}		
		\have [7] {} {\forall x \ \ \mathcal{P}(x)} \by{\text{Introduction de $\forall$}}{}
		\have [8] {} {\lantitauto} \by{\text{Contradiction 1 avec 7}}{}
	\close		
	\have [\lturn] {} {\exists x \ \ \lnot \mathcal{P}(x)} \by{\text{Par contradiction de 2-8}}{}
	\end{nd}
	$
	
\end{exemple}



%%%%%%%%%%%%%%%%%%%%%%%%%%%%%%%%%%%%%%%%%%%%%%%%%%%%%%%%%%%%%%%%%%%%%
\section{Détails des règles}


On reprend les règles une par une avec davantage d'explications et des exemples.

%--------------------------------------------------------------------
\subsection{Règle I. Élimination de $\forall$}


\mybox{$
	\forall x \ \ \mathcal{P}(x)  \quad \lturn \quad  \mathcal{P}(c)
	$}

\medskip

\emph{Conditions.} 
\begin{itemize}
	\item Si $c$ est une constante du langage alors il n'y a pas de condition si la substitution de $x$ par $c$ est valide.
	
	\item Si $c$ est un terme, aucune de ses variables ne doit être liée (capturée) par un quantificateur de $\mathcal{P}$ lors de la substitution.
\end{itemize}

\emph{Autres formulations.}

\[
\forall x \ \ \mathcal{P}(x)  \quad \lturn \quad  \mathcal{P}[c/x]
\]
On trouve aussi parfois l'écriture un peu abusive : $\forall x \ \ \mathcal{P}(x)  \ \lturn \  \mathcal{P}(x)$.

Et voici le schéma en déduction naturelle :
\[
\begin{nd}
	\hypo [~] {} {\forall x \ \ \mathcal{P}(x)}
	\have [\lturn] {} {\mathcal{P}(c)} 
\end{nd}
\]


\medskip

\emph{Exemples.} 
\begin{itemize}

	\item $\forall x \ P(x) \land Q(x,c) \ \ \lturn \ \  P(c) \land Q(c,c)$ est une preuve correcte. On a remplacé les occurrences de $x$ par une constante $c$, même si cette constante apparaissait déjà dans la formule.
	
	\item $\forall x \ P(x)  \ \ \lturn \ \ P( f(a,b) )$ est correct si $f$ est une fonction et $a, b$ deux constantes.
	
	\item $\forall x \ P(x)  \ \ \lturn \ \ P(y)$ est aussi correct, à condition que $y$ ne soit pas capturée par un quantificateur dans $\mathcal{P}$.
\end{itemize}

\medskip

\emph{Contre-exemples.} 
\begin{itemize}
	
	\item $\big( \forall x \ P(x) \big) \limplies Q(a,b)$ ne prouve pas  $P(c) \limplies Q(a,b)$,  car le \og{}pour tout\fg{} n'est pas le connecteur principal de la phrase. 
	Par contre $ \forall x \ \big(P(x) \limplies Q(a,b) \big)$ prouverait bien $P(c) \limplies Q(a,b)$.
	
	\item $\forall x \ \exists y \ Q(x,y)$ ne prouve pas $\exists y \ Q(y,y)$.
	En effet, dans la formule \og{}$\exists y \ Q(x,y)$\fg{}, la variable $x$ apparaît comme variable libre. La substitution de $x$ par $y$ est donc interdite car la variable $y$ injectée lors du remplacement serait capturée par le quantificateur $\exists y$ de la formule.
\end{itemize}

\medskip

\emph{Remarques.} Comment parler d'une constante $c$, si le langage $\mathcal{L}$ n'a pas de constante à ce nom ou bien n'a pas de constante du tout ? On s'autorise à étendre le langage $\mathcal{L}$ en un langage $\mathcal{L}' = \mathcal{L} \cup \{ c \}$ où $c$ est un nouveau symbole de constante.


%--------------------------------------------------------------------
\subsection{Règle II. Introduction de $\forall$}

\mybox{$
	\mathcal{P}(c) \quad \lturn \quad  \forall x \ \ \mathcal{P}(x)
	$}

\medskip

\emph{Conditions.} 
Ce qui est délicat dans de son usage, c'est que $c$ doit être une constante arbitraire, c'est-à-dire une constante dont le symbole n'apparaît dans aucune hypothèse en cours d'utilisation.

	

\emph{Autres formulations.}
	
\[
\mathcal{P}(x) \quad \lturn \quad  \forall x \ \ \mathcal{P}(x)
\]
\[
\begin{nd}
	\hypo [~] {} {\mathcal{P}(c)}
	\have [\lturn] {} {\forall x \ \ \mathcal{P}(x)} 
\end{nd}
\]

Soit $\Gamma$ un ensemble de formules et $c$ une constante qui n'apparaît pas dans les formules de $\Gamma$.
\[
\text{ Si } \ \Gamma \ \ \lturn \ \ \mathcal{P}(c)  \quad 
\text{ alors } \quad \Gamma \ \ \lturn \ \ \forall x \ \ \mathcal{P}(x)
\]

\medskip	
	
\emph{Exemples.}

\begin{itemize}
	\item $c > 0 \ \ \lturn \ \ \forall x \ x > 0$, à condition qu'aucune hypothèse ne  porte sur $c$.
	
	\item $f(a,b) \neq 0  \ \ \lturn \ \ \forall x \ \forall y \ f(x,y) \neq 0$, à condition qu'aucune hypothèse ne porte ni sur $a$, ni sur $b$.
\end{itemize}	


\medskip

\emph{Contre-exemples.}
Beaucoup de choses sont interdites et il est facile de se tromper dans l'application de cette règle.

\begin{itemize}
	\item $c>0, \ P(c)$ ne prouve pas $\forall x \ P(x)$ car il y a la contrainte $c>0$ qui fait que $c$ n'est pas une constante arbitraire.
	Par contre, à partir des hypothèses $c>0$ et $P(c)$, on pourrait prouver que 
	$\forall x \ \big( x > 0 \limplies P(x) \big)$.
	
	\item Voici quasiment le même exemple avec l'écriture en déduction naturelle et un autre exemple :


\medskip
\qquad\qquad
$
\begin{nd}
	\hypo [~] {} {c>0}  
	\have [~] {} {P(c)} 	
	\have [\lturn] {} {\xcancel{\forall x \ \ P(x)}} 
\end{nd}
$
\qquad\qquad
$
\begin{nd}
	\hypo [~] {} {\forall x \ \ Q(c,x)}  
	\have [~] {} {Q(c,c)} 	\by{\text{Élimination de $\forall$ correcte}}{}
	\have [\lturn] {} {\xcancel{\forall x \ \ Q(x,x)}} \by{\text{Introduction de $\forall$ non valide}}{}
\end{nd}
$	

\medskip
	
	\item Il est interdit de ne remplacer que certaines occurrences de $c$. Par exemple, $c=c$ ne prouve pas $\forall x \ x = c$.
\end{itemize}


\medskip
	
\emph{Remarque.}	
Le mieux est que le symbole $c$ utilisé dans la règle n'apparaisse pas du tout auparavant dans la preuve. Mais ce n'est pas toujours aussi simple.
Voici un exemple valide :

\medskip
\qquad\qquad
$
\begin{nd}
	\have [~] {} {\ \vdots} \by{\text{$c$ n'apparaît pas avant}}{}	
	\open		
		\hypo [~] {} {P(c)}  \by{\text{Sous-hypothèse, condition porte sur $c$}}{}
		\have [~] {} {Q(c)} \by{\text{Conclusion dans ce cas}}{}	
	\close 		
	\have [~] {} {P(c) \limplies Q(c)} \by{\text{Il n'y a plus aucune hypothèse en cours d'utilisation contenant $c$}}{}
	\have [\lturn] {} {\forall x \ \ \big( P(x) \limplies Q(x) \big)} \by{\text{Introduction de $\forall$}}{}
\end{nd}
$



%--------------------------------------------------------------------
\subsection{Règle III. Introduction de $\exists$}

\mybox{$
	\mathcal{P}(c) \quad \lturn \quad  \exists x \ \ \mathcal{P}(x)
	$}

\medskip
	
\emph{Conditions.} $c$ est une constante ou bien un terme sans variable libre.
	
\medskip
	
\emph{Autres formulations.}

\[
\mathcal{P}[c/x] \quad \lturn \quad  \exists x \ \ \mathcal{P}(x)
\]	
\[
\begin{nd}
	\hypo [~] {} {\mathcal{P}(c)}
	\have [\lturn] {} {\exists x \ \ \mathcal{P}(x)}
\end{nd}
\]


\medskip

\emph{Exemples.}

\begin{itemize}
	\item $c^2 \neq  c \ \ \lturn \ \ \exists x \ x^2 \neq x$ (déjà vu).
	
	\item $c^2 \neq  c \ \ \lturn \ \ \exists x \ x^2 \neq c$ 
	(remplacement partiel autorisé)
	
	\item $c^2 \neq  c \ \ \lturn \ \ \exists x \ c^2 \neq x$ (remplacement partiel)
	
	\item $a \times b = 12 \ \ \lturn \ \ \exists x \ \exists y \ x \times y  = 12$ (deux utilisations de la règle)
\end{itemize}
 	
\medskip
 
\emph{Remarque.} Il n'y a pas de danger avec cette règle. Il faut cependant prendre soin d'introduire le \og{}il existe\fg{} en tête de la formule obtenue. Par exemple, ceci est correct :
\[
\forall x \ \ Q(x,b)  \quad \lturn \quad  \exists y \ \ \forall x \ \ Q(x,y)
\]
Mais la même prémisse ne prouverait pas $\forall x \ \ \exists y \ \ Q(x,y)$.

\medskip 
\emph{Domaine non vide.} En logique du premier ordre, on supposera toujours qu'on travaille avec un domaine non vide.


%--------------------------------------------------------------------
\subsection{Règle IV. Élimination de $\exists$}

\mybox{$
	\exists x \ \ \mathcal{P}(x), \ \big[\mathcal{P}(c) \lturn \mathcal{Q} \big]  \quad \lturn \quad  \mathcal{Q}
	$}

La notation entre crochets correspond à une sous-preuve.
 	
\medskip


\emph{Conditions.} $c$ doit être une \defi{constante fraîche}\index{constante fraiche@constante fraîche} : d'une part elle ne doit apparaître dans aucune hypothèse en cours d'utilisation, d'autre part elle ne doit pas apparaître dans la formule $\mathcal{Q}$ (ni bien sûr dans $\mathcal{P}(x)$). Le mieux est que la constante n'apparaisse ni avant dans la preuve, ni dans $\mathcal{Q}$.


\emph{Autres formulations.}
\[
\exists x \ \ \mathcal{P}(x), \ \big[\mathcal{P}[c/x] \lturn \mathcal{Q} \big]  \quad \lturn \quad  \mathcal{Q}
\]
\[
\exists x \ \ \mathcal{P}(x), \ \big[\mathcal{P}(x) \lturn \mathcal{Q} \big]  \quad \lturn \quad  \mathcal{Q}
\]

Soit $\Gamma$ un ensemble de formules, $c$ une constante du langage qui n'apparaît ni dans $\mathcal{Q}$, ni dans les formules de $\Gamma$ :
\[
\text{ Si } \quad \Gamma \ \ \lturn \ \exists x \ \ \mathcal{P}(x)
\quad \text{ et } \quad \mathcal{P}(c) \ \ \lturn \ \ \mathcal{Q},  \quad 
\text{ alors } \ \Gamma \ \ \lturn \ \ \mathcal{Q}
\]

Le plus parlant est l'écriture en déduction naturelle.
\[
\begin{nd}
	\hypo [~] {} {\exists x \ \ \mathcal{P}(x)}  \by{}{}
	\open				
	\hypo [~] {} {\mathcal{P}(c)} \by{\text{Sous-hypothèse}}{}
	\have [~] {} {\ \vdots} \by{}{}
	\have [~] {} {\mathcal{Q}} \by{}{}
	\close 
	\have [\lturn] {} {\mathcal{Q}} \by{}{}
\end{nd}
\]

	
\medskip

\emph{Exemple.} Voyons comment utiliser cette règle dans la pratique :
\[
\exists x \ \ f(x)=0, \ \forall x \ \ \big( f(x) = 0 \limplies g(x) = 1 \big) \quad \lturn \quad  \exists x \ \ g(x) = 1
\]
Ce n'est pas une application directe de la règle, mais presque :
 	
\medskip
\qquad\qquad
$
\begin{nd}
	\hypo [~] {} {\exists x \ \ f(x) = 0}  \by{}{}
	\hypo [~] {} {\forall x \ \ \big( f(x) = 0 \limplies g(x) = 1 \big)}  \by{}{}
	\open				
	\hypo [~] {} {f(c) = 0} \by{\text{Sous-hypothèse}}{}
	\have [~] {} {f(c) = 0 \limplies g(c) = 1} \by{\text{Élimination de $\forall$}}{}	
	\have [~] {} {g(c) = 1} \by{\text{Par l'implication}}{}
	\have [~] {} {\exists x \ \ g(x) = 1} \by{\text{C'est la formule $\mathcal{Q}$ de la règle}}{}
	\close 
	\have [\lturn] {} {\exists x \ \ g(x) = 1} \by{\text{Élimination de $\exists$}}{}
\end{nd}
$
 	
\medskip

\emph{Remarque.}
Si on enchaîne plusieurs \og{}il existe\fg{}, on doit prendre une nouvelle constante fraîche pour chacun. 
Par exemple, si on a :
\[
\exists x  \ \ P(x), \ \  \exists y \ \ Q(y)\]
alors dans une preuve on ne peut pas choisir $P(c)$ et $Q(c)$ avec le même $c$. On doit utiliser deux constantes fraîches $c$ et $d$ afin d'écrire $P(c)$ et $Q(d)$ dans les sous-hypothèses.


%--------------------------------------------------------------------
\subsection{Règle V. Égalité}

\mybox{$
	t = s, \ \mathcal{P}(t) \quad \lturn \quad  \mathcal{P}(s)
	$}
	
 	
\medskip

\emph{Conditions.} $t$ et $s$ sont des termes du langage. La substitution de $x$ par $t$ et $s$ doit être valide. Il n'y a pas d'autres conditions.

	
\emph{Autres formulations.}
\[
x = y, \ \mathcal{P}(x) \quad \lturn \quad  \mathcal{P}(y)
\]
\[
\begin{nd}
	\hypo [~] {} {t=s}
	\hypo [~] {} {\mathcal{P}(t)}	
	\have [\lturn] {} {\mathcal{P}(s)} 
\end{nd}
\]
	
\medskip

\emph{Exemple.} $t=x+1, \ t^2 \neq 0 \ \ \lturn \ \ (x+1)^2 \neq 0$.
 	
\medskip

\emph{Remarque.} Cette règle peut être remplacée par les axiomes suivants pour le symbole égalité :
\begin{itemize}
	\item $\forall x \ \ x = x$
	
	\item $\forall x \ \ \forall y \ \ \big( (x = y) \limplies (\mathcal{P}(x) \lequiv \mathcal{P}(y)) \big)$
\end{itemize}



%%%%%%%%%%%%%%%%%%%%%%%%%%%%%%%%%%%%%%%%%%%%%%%%%%%%%%%%%%%%%%%%%%%%%
\section*{Exercices}

\subsection*{Preuves en logique du premier ordre}

\begin{exercicecours}
Écrire les preuves des théorèmes suivants. $P, Q, R$ désignent des prédicats, $f, g$ des fonctions.
\begin{enumerate}
	
	\item $\forall x \ P(x), \ \ P(a) \limplies Q(a) \ \ \lturn \ \ Q(a)$ 
	
	\item $\forall x \ P(x) \land Q(x) \ \ \lturn \ \ \exists x \ P(x)$
	
	\item $\forall x \ \forall y \ R(x,y) \ \ \lturn \ \ \forall x \ R(x,x)$
	
	\item $\exists x  \ f(x)=g(x) \ \ \lturn \ \ \exists x \ \exists y \ f(x)=g(y)$

\end{enumerate}
\end{exercicecours}


\begin{exercicecours}
	Écrire les preuves des théorèmes suivants. $P, Q, R$ désignent des prédicats, $f$ une fonction.
	\begin{enumerate}
		
		\item $\lturn \forall x \ \ \forall y \ \ \big( f(x) \neq f(y) \  \limplies \ x \neq y \big)$
		
		\item $\forall x \ P(x), \ \forall y \ Q(y) \ \ \lturn \ \ \exists z \ P(z) \land Q(z)$
		
		\item $\forall x \ P(x) \limplies Q(x), \ \ \forall y \ Q(y) \limplies R(y)
		\ \ \lturn \ \ \forall z \ P(z) \limplies R(z)$
		
		\item $\forall x \ \big( P(x) \limplies Q \big), \ \ \exists x  \ P(x) \ \ \lturn \ \ Q$, où $Q$ ne dépend pas de $x$.
		
	\end{enumerate}
\end{exercicecours}


\subsection*{Détails des règles}


\begin{exercicecours}
Expliquer l'erreur faite en appliquant la règle d'élimination de $\forall$.
\begin{enumerate}
	\item $\exists x \ \forall y \ Q(x,y)$ prouve directement $\exists x \ Q(x,x)$. Bonus : donner une preuve correcte.
	
	\item $\forall x \ \big( P(x) \limplies Q(c) \big)$ prouve $P(c)$.
	
	\item $\forall x \ \exists y \  Q(x,y)$ prouve $\exists y \  Q(y,y)$.
\end{enumerate}
\end{exercicecours}


\begin{exercicecours}
Expliquer l'erreur faite	en appliquant la règle d'introduction de $\forall$.
\begin{enumerate}
	\item $c=1, \ \ c > 0$ prouve $\forall x \ x> 0$. 
	
	\item $\exists x \ P(x)$ prouve $\forall x \ P(x)$. 
	
	\item $(x \neq 0) \limplies (x^2 \neq 0)$ prouve $\forall x \ x^2 \neq 0$. 
	
	\item $(x = 0) \lor (x \neq 0)$ prouve $\big( \forall x \ x = 0 \big) \lor \big( \forall x \ x \neq 0 \big)$. 		
	
	\item $f(c)=g(c)$ prouve $\forall x \ f(x)=g(c)$. 
\end{enumerate}
\end{exercicecours}


\begin{exercicecours}
Expliquer l'erreur faite	en appliquant la règle d'introduction de $\exists$.
\begin{enumerate}
	\item $\forall x \ P(x)$ prouve directement $\exists x \ P(x)$. 
	
	\item $P(a) \land Q(a)$ prouve $\exists x \ P(x) \land Q(y)$.
	
	\item $c=0, \ \ \exists x \ P(x)$ prouve $\exists x \ (x=0) \land P(x)$.
	
	\item $\forall x \ Q(y,y)$ prouve $\exists x \ \forall y \ Q(x,y)$.
	
	\item $\forall x \ Q(y,y)$ prouve $\forall y \ \exists x \ Q(x,y)$.		
	
\end{enumerate}
\end{exercicecours}


\begin{exercicecours}
Expliquer l'erreur faite	en appliquant la règle d'élimination de $\exists$.
\begin{enumerate}
	\item $\exists x \ P(x), \ \ P(0) \limplies Q$ prouve $Q$. 
	
	\item $\exists x \ P(x)$ prouve $P(c)$. 
	
	\item $\exists x \ \big( P(x) \limplies Q \big), \ \ P(c)$ prouve $Q$.
	
	\item $\exists x \ P(x), \ \ \exists y \ Q(y)$ prouve $\exists z \ P(z) \land Q(z)$. 	 		
\end{enumerate}
\end{exercicecours}


\subsection*{Négation des quantificateurs et règles dérivées}


\begin{exercicecours}
Écrire ce que prouvent les formules suivantes. 
La formule prouvée ne doit plus contenir aucune négation à gauche d'un quantificateur.
\begin{enumerate}
	\item $\lnot \Big( \forall x \ f(x) = g(x) \Big)$
	
	\item $\lnot \Big( \exists x \ P(x) \lequiv Q(x) \Big)$
	
	\item $\lnot \big[ \forall x \ \big( P(x) \limplies \exists y \ Q(y) \big) \big]$
\end{enumerate}
\end{exercicecours}


\begin{exercicecours}
Prouver toutes les règles dérivées à l'aide des règles de base.	
\end{exercicecours}

\end{document}
